\section{Motivation}

Multi-agent systems (MAS) have vast potential for various applications of robotics.  From search and rescue \cite{drew_multi-agent_2021}, to warehouse logistics \cite{dusadeerungsikul_multi-agent_2022} MAS can be used to improve efficiency and effectiveness in various domains.  However, one challenge with MAS is the significant overhead in deployment and developing such systems.

Because deployment of MAS is often a complex and expensive task, most of the body of research in MAS is completed in simulation.  Specifically, a large body of research exists in real-time strategy games (specifically StarCraft) \cite{jensen_andreas_schmidt_interfacing_2015}.  Within this body of research, the game defines a set of constraints and actions that agents can take while providing reinforcement learning algorithms significant feedback on the performance of the agents \cite{zhou_hierarchical_2021}.  Because of the success of training agents in games, there has become a growing interest in applying these techniques to more realistic and increasingly complex scenarios \cite{goecks_games_2023}.  However, training agents in games is still complex and expensive and many real-world deployments require a human in the loop to ensure that the agents are performing as expected.  This is especially true in scenarios where the agents are operating in a dynamic environment with changing constraints and objectives.

Branching out of the need for a human in the loop, there has been a growing interest in using large language models (LLMs) to control agents in MAS.  In most research, this LLM control has taken the form of high level task allocation, mid level motion planning, and low level action generation \cite{li_large_2026}.  For this project, we will be focusing on the high level task allocation problem.  This problem has increasingly been shown to be accomplishable by LLMs, with LLMs being able to generate task allocation plans for MAS in a variety of standard MAS and swarm problem formulations (such as flocking and formation control) \cite{wang_decentralized_2026}, \cite{eumi_swarmchat_2026}, \cite{ji_genswarm_2026} as well as robot soccer \cite{brienza_llcoach_2025}.  Although many different approaches have been attempted for LLM control of MAS, a particularly interesting approach involves having LLMs generate plans using a sort of intermediate representation.  In \cite{majid_commandswarm_2026}, the authors have the LLM generate a behavior tree for the agents to follow.  In \cite{xu_nl2hltl2plan_2025}, the authors use Linear Temporal Logic (LTL) as an intermediate representation for the LLM to generate plans.

While LLMs have been shown to be able to generate plans for MAS, there has been little work done in evaluating these systems with a human in the loop and in situations with many robots.  First, much of the previous work has focused on either pre-defined and common MAS problems (such as flocking and formation control) \cite{ji_genswarm_2026} or on a small number of robots without a human in the loop \cite{wang_decentralized_2026}, \cite{eumi_swarmchat_2026}, \cite{borate_comuros_2026}.  Second, there has been little work done in evaluating the user-side metrics of these systems.  In \cite{xu_nl2hltl2plan_2025}, the authors conducted a 4 participant user study to evaluate the usability of their system.  However, this study was limited to the performance of the system and did not evaluate the user experience of the participants.

Therefore, in this project, we will be evaluating and comparing the performance and user-experience of a traditional real-time strategy game interface and a LLM interface for controlling a MAS in a simulated game environment.  Specifically, we are interested in investigating whether utilizing an LLM interface for controlling the MAS leads to decreased task load and what impact the LLM interface has on performance.  Additionally, much of the previous research on using LLMs to control MAS has focused on contrived examples so we are additionally interested in investigating how humans interact and perceive the LLM interface when controlling a MAS in a more complex scenario.
